% !TeX root = part5.tex
\documentclass[../final.tex]{subfiles}

\begin{document}

% Лекция 29.09.2026. Продолжение главы об интерполяции.
\ifSubfilesClassLoaded{\setcounter{chapter}{4}\setcounter{section}{6}}{}
\captionsetup{font=footnotesize,labelfont=bf,skip=5pt}

\subsection{Узлы Чебышёва}

Остаточный член интерполяции содержит \rconcept{узловой множитель}
\begin{equation*}
    \omega_{n+1}(x)=\prod_{i=0}^{n}(x-x_i),\qquad
    y(x)-P_n(x)=\frac{y^{(n+1)}(\xi)}{(n+1)!}\omega_{n+1}(x).
\end{equation*}
Поэтому при выборе узлов стремятся минимизировать величину
\begin{equation*}
    \max_{x\in[a,b]}\left|\prod_{i=0}^{n}(x-x_i)\right|.
\end{equation*}
На отрезке $[-1,1]$ минимум достигается, если $x_i$ --- корни
полинома Чебышёва $T_{n+1}$:
\begin{equation*}
    T_{n+1}(t)=\cos\bigl((n+1)\arccos t\bigr),\qquad
    t_i=\cos\frac{(2i+1)\pi}{2(n+1)},\quad i=0,\ldots,n.
\end{equation*}
Эти узлы сгущаются у концов отрезка.
На произвольном отрезке $[a,b]$ упорядочим их по возрастанию:
\begin{equation*}
    x_i=\frac{a+b}{2}-\frac{b-a}{2}
    \cos\frac{(2i+1)\pi}{2(n+1)},\qquad i=0,\ldots,n.
\end{equation*}
Для такого выбора
\begin{equation*}
    \max_{x\in[a,b]}|\omega_{n+1}(x)|
    =\frac{(b-a)^{n+1}}{2^{2n+1}}.
\end{equation*}
Таким образом, \rresult{узлы Чебышёва минимизируют максимум модуля
узлового множителя} и уменьшают соответствующую оценку погрешности.

\section{Численное дифференцирование}

\subsection{Производные интерполяционного полинома}

Пусть известны значения $y(x_i)=y_i$ в узлах $x_0,\ldots,x_n$.
Вспомним полином Ньютона:
\begin{equation*}
    P_n(x)=y(x_0)+\sum_{k=1}^{n}
    y^{(k)}(x_0,\ldots,x_k)\prod_{j=0}^{k-1}(x-x_j).
\end{equation*}
При нескольких аргументах $y^{(k)}(x_0,\ldots,x_k)$ означает
разделённую разность, а $y^{(k)}(x)$ с одним аргументом --- производную.
Приближённые значения производных получаем дифференцированием $P_n$.

Введём \rassumption{локальную равномерную сетку с шагом $h=x_{i+1}-x_i$}.
При фиксированном числе близких узлов и ограниченных производных
разложение по формуле Тейлора даёт
\begin{equation*}
    y(x)=P_n(x)+O(h^{n+1}),\qquad
    y^{(m)}(x)=P_n^{(m)}(x)+O(h^{n+1-m}),\quad 1\le m\le n.
\end{equation*}
Для полинома степени $m$ по $m+1$ узлам старший коэффициент равен
разделённой разности порядка $m$, поэтому
\begin{align*}
    P_m^{(m)}(x)&=m!\,y^{(m)}(x_0,\ldots,x_m),\\
    y^{(m)}(x)&=m!\,y^{(m)}(x_0,\ldots,x_m)+O(h).
\end{align*}

\subsection{Первая производная: правая и левая разности}

Правая разностная производная определяется как
\begin{equation*}
    \Delta_{\mathrm r}y(x)
    =\frac{y(x+h)-y(x)}{h}.
\end{equation*}
Это наклон секущей через точки с абсциссами $x$ и $x+h$.
Аналогично \rconcept{левая разностная производная} равна
\begin{equation*}
    \Delta_{\mathrm l}y(x)
    =\frac{y(x)-y(x-h)}{h}.
\end{equation*}

\begin{rbox}[breakable=false,left=18pt,right=18pt,top=14pt,bottom=10pt]{[]}
    \centering
    \begin{tikzpicture}[x=1.75cm,y=0.85cm,
        every node/.style={text=rtext,font=\small}]
        \draw[rtext,-{Stealth[length=2.2mm]}] (0,0)--(4.0,0) node[right] {$x$};
        \draw[rtext,-{Stealth[length=2.2mm]}] (0,0)--(0,2.7) node[left] {$y$};
        \draw[rc4,line width=1.4pt,domain=0.3:3.6,samples=90]
            plot(\x,{0.45+0.15*\x*\x});
        \draw[p3,line width=1.1pt] (1,0.6)--(2,1.05);
        \draw[red3,line width=1.1pt] (2,1.05)--(3,1.8);
        \foreach \xx/\yy/\lab in {1/0.6/{x-h},2/1.05/{x},3/1.8/{x+h}}{
            \draw[rtext!45,dashed] (\xx,0)--(\xx,\yy);
            \fill[rtext] (\xx,\yy) circle(1.6pt);
            \node[below=4pt] at (\xx,0) {$\lab$};
        }
        \node[p3,above left] at (1.5,0.825) {$\Delta_{\mathrm l}y$};
        \node[red3,above left] at (2.6,1.5) {$\Delta_{\mathrm r}y$};
    \end{tikzpicture}
    \captionof{figure}{Левая и правая разностные производные --- наклоны
    секущих на соседних участках сетки.}
    \label{fig:ovf5_one_sided_differences}
\end{rbox}

\subsection{Оценка погрешности первой производной}

Погрешность правой разности равна $R_{\mathrm r}(h)=
\Delta_{\mathrm r}y(x)-y'(x)$. Разложим функцию в ряд Тейлора:
\begin{equation*}
    y(x+h)=y(x)+hy'(x)+\frac{h^2}{2}y''(x)
    +\frac{h^3}{6}y'''(x)+O(h^4).
\end{equation*}
Подстановка даёт
\begin{equation*}
    R_{\mathrm r}(h)
    =\frac h2y''(x)+\frac{h^2}{6}y'''(x)+O(h^3)=O(h).
\end{equation*}
Для левой разности аналогично
\begin{equation*}
    R_{\mathrm l}(h)=\Delta_{\mathrm l}y(x)-y'(x)
    =-\frac h2y''(x)+\frac{h^2}{6}y'''(x)+O(h^3)=O(h).
\end{equation*}
Обе односторонние формулы имеют \rresult{первый порядок точности}.

\subsection{Центральная разность}

Усреднение левой и правой разностей устраняет член первого порядка:
\begin{equation*}
    \Delta_{\mathrm c}y(x)
    =\frac{\Delta_{\mathrm r}y(x)+\Delta_{\mathrm l}y(x)}2
    =\frac{y(x+h)-y(x-h)}{2h}.
\end{equation*}
Для достаточно гладкой функции погрешность равна
\begin{equation*}
    R_{\mathrm c}(h)=\Delta_{\mathrm c}y(x)-y'(x)
    =\frac{h^2}{6}y'''(x)+O(h^4)=O(h^2).
\end{equation*}
Центральная разность имеет \rresult{второй порядок точности}.

\subsection{Вторая производная}

По трём произвольным близким узлам
\begin{equation*}
    y''(x)=2!\,y^{(2)}(x_0,x_1,x_2)+O(h).
\end{equation*}
На симметричной равномерной сетке $x_0=x-h$, $x_1=x$, $x_2=x+h$
разделённая разность равна
\begin{align*}
    y^{(2)}(x-h,x,x+h)
    &=\frac{y^{(1)}(x-h,x)-y^{(1)}(x,x+h)}{(x-h)-(x+h)}\\
    &=\frac{y(x+h)-2y(x)+y(x-h)}{2h^2}.
\end{align*}
Центральная формула и её разложение по Тейлору имеют вид
\begin{align*}
    \Delta_2y(x)&=\frac{y(x+h)-2y(x)+y(x-h)}{h^2},\\
    \Delta_2y(x)&=y''(x)+\frac{h^2}{12}y^{(4)}(x)+O(h^4).
\end{align*}
Благодаря симметрии сетки эта формула также имеет
\rresult{второй порядок точности}.

\chapter{Численное решение обыкновенных дифференциальных уравнений}
\rsubtitle{Метод Эйлера и явный метод средней точки}

\section{Задача Коши и равномерная сетка}

Рассмотрим \rconcept{задачу Коши} для ОДУ первого порядка:
\begin{equation*}
    \begin{cases}
        u'(x)=f(x,u(x)),\qquad x\in[a,b],\\
        u(a)=u_0.
    \end{cases}
\end{equation*}
Функцию $u(x)$ будем считать точным решением.
Значения численного решения на сетке обозначим через $y_i$.
Для $N$ равных шагов
\begin{equation*}
    h=\frac{b-a}{N},\qquad x_i=a+ih,\qquad i=0,\ldots,N,
    \qquad y_i\simeq u(x_i).
\end{equation*}

\subsection{Производные точного решения}

Если $f$ достаточно гладкая, производные $u$ можно выразить через $f$.
По правилу дифференцирования сложной функции
\begin{align*}
    u'&=f,\\
    u''&=f_x+f_u u'=f_x+ff_u,\\
    u'''&=f_{xx}+2ff_{xu}+f^2f_{uu}+f_xf_u+f(f_u)^2.
\end{align*}
Все значения $f$ и её частных производных в этих формулах берутся
в точке $(x,u(x))$.

\section{Метод Эйлера}

\subsection{Разностная схема}

Заменим производную правой разностью и вычислим правую часть
в начале шага:
\begin{equation*}
    \frac{y_{i+1}-y_i}{h}=f(x_i,y_i).
\end{equation*}
Получаем \rconcept{явную схему Эйлера}:
\begin{equation*}
    \begin{cases}
        y_{i+1}=y_i+h f(x_i,y_i),\qquad i=0,\ldots,N-1,\\
        y_0=u_0.
    \end{cases}
\end{equation*}
На каждом шаге наклон решения заменяется наклоном касательной
в начале этого шага.

\begin{rbox}[breakable=false,left=18pt,right=18pt,top=14pt,bottom=10pt]{[]}
    \centering
    \begin{tikzpicture}[x=1.4cm,y=0.8cm,
        every node/.style={text=rtext,font=\small}]
        \draw[rtext,-{Stealth[length=2.2mm]}] (0,0)--(5.9,0) node[right] {$x$};
        \draw[rtext,-{Stealth[length=2.2mm]}] (0,0)--(0,4.55) node[left] {$u$};
        \draw[rc4,line width=1.4pt,domain=0.6:5.25,samples=90]
            plot(\x,{0.4+0.14*\x*\x});
        \draw[p3,line width=1.3pt] (2.2,1.0776)--(3.1,1.632);
        \draw[p3,dashed] (3.1,1.632)--(3.8,2.0632);
        \draw[rtext!45,dashed] (0,1.0776)--(2.2,1.0776);
        \node[left] at (0,1.0776) {$y_i$};
        \draw[rtext!45,dashed] (2.2,0)--(2.2,1.0776);
        \draw[rtext!45,dashed] (3.1,0)--(3.1,1.7454);
        \fill[rtext] (2.2,1.0776) circle(1.6pt);
        \fill[p3] (3.1,1.632) circle(1.6pt);
        \fill[rc4] (3.1,1.7454) circle(1.6pt);
        \node[below=4pt] at (0.8,0) {$a$};
        \node[below=4pt] at (2.2,0) {$x_i$};
        \node[below=4pt] at (3.1,0) {$x_{i+1}$};
        \node[below=4pt] at (5.2,0) {$b$};
        \node[rc4,above left] at (4.5,3.235) {$u(x)$};
        \node[p3,right] at (3.8,2.0632) {шаг Эйлера};
    \end{tikzpicture}
    \captionof{figure}{Шаг метода Эйлера: переход вдоль касательной
    вместо движения по точной интегральной кривой.}
    \label{fig:ovf5_euler_step}
\end{rbox}

\subsection{Локальная погрешность}

Разложим точное решение около $x_i$:
\begin{equation*}
    u(x_i+h)=u(x_i)+hu'(x_i)+\frac{h^2}{2}u''(x_i)+O(h^3).
\end{equation*}
Пусть в начале одного шага $y_i=u(x_i)$. Тогда
\begin{equation*}
    y_{i+1}=u(x_i)+h f(x_i,u(x_i))=u(x_i)+hu'(x_i).
\end{equation*}
Погрешность одного шага равна
\begin{equation*}
    u(x_{i+1})-y_{i+1}=\frac{h^2}{2}u''(x_i)+O(h^3).
\end{equation*}
Таким образом, \rconcept{локальная погрешность} метода Эйлера
имеет порядок $O(h^2)$.

\subsection{Невязка и порядок аппроксимации}

Введём разностный оператор
\begin{equation*}
    \Delta_h v_i=\frac{v_{i+1}-v_i}{h}-f(x_i,v_i).
\end{equation*}
Численная схема записывается как $\Delta_h y_i=0$.
Подстановка точного решения $u_i=u(x_i)$ в этот оператор даёт
\rconcept{невязку}:
\begin{equation*}
    \Psi_{h,i}=\Delta_h u_i
    =\frac{u(x_i+h)-u(x_i)}h-f(x_i,u(x_i)).
\end{equation*}
Невязка численной аппроксимации дифференциального уравнения ---
результат подстановки его точного решения в разностную схему.
Если $\Psi_{h,i}=O(h^p)$, то схема аппроксимирует уравнение
с порядком $p$.

Для метода Эйлера
\begin{align*}
    \Psi_{h,i}
    &=u'(x_i)+\frac h2u''(x_i)+O(h^2)-f(x_i,u(x_i))\\
    &=\frac h2u''(x_i)+O(h^2)=O(h).
\end{align*}
Следовательно, схема Эйлера имеет \rresult{первый порядок аппроксимации}.
Невязка отличается от погрешности одного шага множителем $1/h$.

\subsection{Накопление погрешности}

На фиксированном отрезке число шагов равно $N=(b-a)/h$.
Для оценки накопления локальных ошибок необходимо учитывать
их распространение между шагами.
Предположим, что $f$ удовлетворяет \rassumption{условию Липшица по $u$}:
\begin{equation*}
    |f(x,v)-f(x,w)|\le L|v-w|.
\end{equation*}
Если $e_i=y_i-u(x_i)$ и локальная погрешность ограничена величиной $Ch^2$,
то
\begin{equation*}
    |e_{i+1}|\le(1+Lh)|e_i|+Ch^2,\qquad e_0=0.
\end{equation*}
При $L>0$ отсюда следует
\begin{equation*}
    |e_i|\le\frac{Ch}{L}\bigl((1+Lh)^i-1\bigr)
    \le\frac{Ch}{L}\bigl(e^{L(b-a)}-1\bigr)=O(h).
\end{equation*}
При $L=0$ оценка принимает вид
\begin{equation*}
    |e_N|\le NCh^2=C(b-a)h=O(h).
\end{equation*}
Таким образом, глобальная погрешность Эйлера на фиксированном отрезке
имеет \rresult{первый порядок по шагу $h$}.

\section{Модифицированный метод Эйлера}

\subsection{Интегральная форма уравнения}

Интегрируя $u'=f(x,u)$ на одном шаге, получаем
\begin{equation*}
    u(x_{i+1})=u(x_i)+\int_{x_i}^{x_{i+1}}f(x,u(x))\,dx.
\end{equation*}
Замена интеграла формулой левых прямоугольников даёт
\begin{equation*}
    u(x_{i+1})=u(x_i)+h f(x_i,u(x_i))+O(h^2),
\end{equation*}
то есть обычную схему Эйлера.

Чтобы повысить порядок точности, используем \rconcept{метод средних
прямоугольников}:
\begin{equation*}
    u(x_{i+1})=u(x_i)
    +h f\!\left(x_i+\frac h2,u\!\left(x_i+\frac h2\right)\right)
    +O(h^3).
\end{equation*}
Значение решения в середине шага заранее неизвестно. Оценим его
полушагом метода Эйлера:
\begin{equation*}
    y_{i+1/2}=y_i+\frac h2 f(x_i,y_i).
\end{equation*}

\subsection{Явная схема средней точки}

Получаем \rconcept{модифицированный метод Эйлера}:
\begin{equation*}
    \begin{cases}
        y_{i+1/2}=y_i+\dfrac h2f(x_i,y_i),\\[0.5em]
        y_{i+1}=y_i+h f\!\left(x_i+\dfrac h2,y_{i+1/2}\right),\\[0.5em]
        y_0=u_0.
    \end{cases}
\end{equation*}
В одной формуле
\begin{equation*}
    y_{i+1}=y_i+h f\!\left(x_i+\frac h2,
    y_i+\frac h2 f(x_i,y_i)\right).
\end{equation*}
Здесь правая часть вычисляется дважды: в начале шага и в приближённой
средней точке.

\subsection{Порядок точности}

Если $y_i=u(x_i)$, то разложение $f$ в средней точке даёт
\begin{equation*}
    f\!\left(x_i+\frac h2,u(x_i)+\frac h2f(x_i,u(x_i))\right)
    =f+\frac h2(f_x+ff_u)+O(h^2).
\end{equation*}
Следовательно,
\begin{equation*}
    y_{i+1}=u(x_i)+hu'(x_i)+\frac{h^2}{2}u''(x_i)+O(h^3).
\end{equation*}
Первые три члена совпадают с разложением точного решения.
Поэтому локальная погрешность равна $O(h^3)$, невязка --- $O(h^2)$,
а глобальная погрешность при тех же условиях гладкости и Липшица --- $O(h^2)$.
Модифицированный метод Эйлера имеет \rresult{второй порядок точности}.

\end{document}
